\section{Non-Injective Projection and the Failure of Bell Factorization}
  \label{sec:non-injective-projection}

\section{Non-Injective Projection and the Failure of Factorization}
  \label{sec:noninjective-projection}

  We consider a minimal abstract setting in which observable descriptions arise from
  an underlying configuration space through a generally many-to-one mapping.
  Let $\Omega$ denote the space of underlying configurations, assumed to be fully
  specified and governed by local dynamics.
  Observable outcomes are defined through a projection
  \begin{equation}
    \Pi : \Omega \rightarrow \mathcal{O},
  \end{equation}
  where $\mathcal{O}$ is the space of observable states or measurement outcomes.
  Crucially, $\Pi$ is assumed to be non-injective, so that multiple distinct elements
  of $\Omega$ correspond to the same observable description.

  Let $A$ and $B$ denote two spacelike separated measurement settings, with corresponding
  observable outcomes $a \in \mathcal{O}_A$ and $b \in \mathcal{O}_B$.
  In standard local hidden-variable models, one assumes the existence of a variable
  $\lambda$ such that the joint probability distribution factorizes as
  \begin{equation}
    P(a,b \mid A,B,\lambda) = P(a \mid A,\lambda)\,P(b \mid B,\lambda).
  \end{equation}
  Bell inequalities follow from this factorization together with statistical independence
  assumptions on $\lambda$.

  In the present framework, no such factorization is generically available.
  Because the observable outcomes are obtained through the projection $\Pi$, the
  probability distribution accessible at the effective level is given by
  \begin{equation}
    P(a,b \mid A,B) =
    \sum_{\omega \in \Pi^{-1}(a,b)}
    \mu(\omega),
  \end{equation}
  where $\mu$ is a probability measure on $\Omega$, and $\Pi^{-1}(a,b)$ denotes the
  set of underlying configurations compatible with the joint observable outcome $(a,b)$.

  When $\Pi$ is non-injective, the preimage $\Pi^{-1}(a,b)$ does not decompose into
  independent subsets associated with $A$ and $B$.
  As a result, the joint probability $P(a,b \mid A,B)$ does not admit a representation
  as a product of marginal probabilities conditioned on any set of effective variables.
  The failure of factorization arises from the structure of the projection itself,
  not from any dynamical interaction between the measurement regions.

  Importantly, locality at the underlying level is preserved.
  The measure $\mu(\omega)$ may be defined over configurations with strictly local
  causal structure, and no superluminal influence or signal exchange is required.
  The violation of Bell inequalities therefore reflects a mismatch between the
  factorizable structure assumed in Bell-type derivations and the actual structure
  of effective descriptions induced by non-injective mappings.

  This observation identifies non-injectivity as a sufficient condition for the
  failure of Bell factorization.
  Bell inequalities constrain only those effective models in which the mapping
  from underlying configurations to observables is injective or admits an equivalent
  factorizable representation.
  When this condition fails, Bell-type correlations may arise without violating
  local causality at the underlying level.
